%!TEX program = xelatex
% 本文件通过 run_all.sh 编译（图片在仓库中为 .png，编译时转为 .jpg）
\documentclass[12pt,a4paper]{ctexart}

\usepackage[margin=2.5cm]{geometry}
\usepackage{textcomp}
\usepackage{amsmath,amssymb,booktabs,graphicx,hyperref,caption,enumitem,array,multirow}
\hypersetup{colorlinks=true,linkcolor=blue,urlcolor=blue,citecolor=blue}

\title{\textbf{旋转圆盘电极动力学参数测定\\—— K$_4$Fe(CN)$_6$/K$_3$Fe(CN)$_6$ 体系}}
\author{}
\date{\today}

\begin{document}
\maketitle
\thispagestyle{empty}

\begin{center}
\small
\vspace{1cm}
\textbf{仪器：} CHI660E 电化学工作站 \quad
\textbf{技术：} Tafel Plot (旋转圆盘电极) \\
\textbf{工作电极：} Pt 盘电极 ($\Phi = 5$ mm) \quad
\textbf{参比电极：} Hg/HgO (1M KCl) \\
\textbf{分析工具：} Python (NumPy, SciPy, Pandas, Matplotlib)
\end{center}

\vspace{1.5cm}
\tableofcontents
\newpage

% ============================================================
\section{摘要}

本报告基于 CHI660E 电化学工作站的旋转圆盘电极（RDE）数据，
对 K$_4$Fe(CN)$_6$/K$_3$Fe(CN)$_6$ 氧化还原探针体系进行系统的动力学参数测定。
实验由 4 组独立操作者分别完成，采用 Pt 盘电极（$\Phi = 5$ mm，
$A = 0.196$ cm$^2$），1M KCl 为支持电解质。

通过 Levich 方程（$|i_l|$ vs $\omega^{1/2}$）测定扩散系数：
$D_{\text{O}}(\text{Fe(CN)}_6^{3-}) = 6.55 \times 10^{-6}$ cm$^2$/s，
$D_{\text{R}}(\text{Fe(CN)}_6^{4-}) = 4.15 \times 10^{-6}$ cm$^2$/s，
与文献值（$7.6\times 10^{-6}$、$6.3\times 10^{-6}$ cm$^2$/s）数量级一致。
Koutecky-Levich 分析结合 Tafel 外推获得的传递系数
$\alpha \approx 0.72$--$0.98$ 偏离理想可逆值（$\alpha = 0.5$），
归因于准可逆体系中低过电位区反反应的贡献。
4 组独立测量在 CHI660E 精度范围内高度一致，
验证了实验方案的可复现性。

\newpage

% ============================================================
\section{引言}

旋转圆盘电极（Rotating Disk Electrode, RDE）是电化学动力学研究的经典技术。
当电极以恒定角速度 $\omega$ 旋转时，电极表面的流体动力学行为可用 Levich 理论
精确描述\cite{bard2001}，使得扩散层厚度可控，从而将扩散贡献与动力学贡献分离。

K$_4$Fe(CN)$_6$/K$_3$Fe(CN)$_6$ 氧化还原电对（Fe(CN)$_6^{3-} + e^- \rightleftharpoons$ Fe(CN)$_6^{4-}$）
是电化学研究中最常用的模型体系之一。该体系具有以下优点：
（1）单电子转移（$n=1$），反应机理简单；
（2）在 Pt 电极上接近可逆；
（3）扩散系数和标准速率常数已有大量文献值可供参照。

本实验参照《电化学实验》教材（唐安平主编，实验 26），
利用 RDE 技术测定该体系的扩散系数 $D_{\text{O}}$、$D_{\text{R}}$，
并通过 Koutecky-Levich 和 Tafel 分析评估动力学参数。

\newpage

% ============================================================
\section{实验部分}

\subsection{试剂与仪器}

\begin{itemize}[leftmargin=1.5cm]
  \item 电解液：0.01 mol/L K$_4$Fe(CN)$_6$ + 0.01 mol/L K$_3$Fe(CN)$_6$ + 1 mol/L KCl
  \item 工作电极：Pt 盘电极，$\Phi = 5$ mm，$A = 0.196$ cm$^2$
  \item 参比电极：Hg/HgO（1M KCl）
  \item 对电极：Pt 丝
  \item 仪器：CHI660E 电化学工作站，Tafel Plot 技术
\end{itemize}

\subsection{测试条件}

\begin{table}[h]
\centering
\caption{电化学测试参数}
\begin{tabular}{ll}
\toprule
参数 & 数值 \\
\midrule
技术 & Tafel Plot \\
电位范围 & $-0.1 \rightarrow 0.5$ V \\
扫描速率 & 0.01 V/s \\
转速 (Part A) & 500, 1000, 1500, 1800, 2000, 2500 rpm \\
浓度 (Part B) & 0.0002--0.01 mol/L @ 2000 rpm \\
温度 & $25 \pm 1$\textdegree C \\
\bottomrule
\end{tabular}
\end{table}

\subsection{实验设计}

实验由 4 组独立操作者分别完成，每组使用指定浓度的 Fe(CN)$_6^{3-/4-}$ 溶液
和 Pt 盘电极独立测试。分析含以下部分：

\begin{enumerate}[leftmargin=1.5cm]
  \item \textbf{Part A -- Levich 分析}：0.01M 浓度，6 个转速 → $D_{\text{O}}$, $D_{\text{R}}$
  \item \textbf{Part B -- 浓度校准}：5 个浓度（0.0002--0.01M），2000 rpm → $|i_l| \propto c^*$
  \item \textbf{Part C -- K-L/Tafel 分析}：0.01M 数据 → K-L 提取 $i_k$ → Tafel → $i_0$, $\alpha$
  \item \textbf{Part D -- 半波电位}：0.0008M 无 KCl → $E_{1/2}$
\end{enumerate}

\newpage

% ============================================================
\section{分析方法}

分析采用 Python 脚本自动化处理，各步骤通过 CSV 文件传递数据：

\begin{enumerate}[leftmargin=1.5cm]
  \item \textbf{数据解析} (\texttt{parse\_chi660e.py})：识别 CHI660E Tafel Plot 格式（3 列数据），
  按实验条件分类汇总
  \item \textbf{Levich 分析} (\texttt{levich\_analysis.py})：提取阴阳极极限电流 $\rightarrow$
  $|i_l|$ vs $\omega^{1/2}$ 线性拟合 $\rightarrow$ 扩散系数 $D_{\text{O}}$、$D_{\text{R}}$
  \item \textbf{浓度校准 + K-L/Tafel} (\texttt{kl\_analysis.py})：
  浓度校准 $\rightarrow$ 选定过电位 $\rightarrow$ K-L 提取 $i_k$
  $\rightarrow$ $\eta$ vs $\log|i_k|$ Tafel 外推 $\rightarrow$ $\alpha$、$i_0$
  \item \textbf{半波电位}：0.0008M 无 KCl 极化曲线 $\rightarrow$ $E_0$、$E_{1/2}$
  \item \textbf{图表} (\texttt{generate\_figures.py})：Matplotlib 生成 Levich/K-L/Tafel/半波电位图
\end{enumerate}

\newpage

% ============================================================
\section{理论背景}

\subsection{Levich 方程}

RDE 上的极限扩散电流由 Levich 方程描述\cite{bard2001}：

\begin{equation}
|i_l| = 0.62\, n\, F\, A\, D^{2/3}\, \nu^{-1/6}\, c^*\, \omega^{1/2}
\label{eq:levich}
\end{equation}

以 $|i_l|$ 对 $\omega^{1/2}$ 作图（Levich 图），斜率为 $0.62 n F A D^{2/3} \nu^{-1/6} c^*$，
由斜率可求解扩散系数 $D$。

\subsection{Koutecky-Levich 方程}

在混合控制区（动力学+扩散），总电流与动力学电流满足：

\begin{equation}
\frac{1}{|i|} = \frac{1}{i_k} + \frac{1}{0.62\, n\, F\, A\, D^{2/3}\, \nu^{-1/6}\, c^*\, \omega^{1/2}}
\label{eq:kl}
\end{equation}

以 $1/|i|$ 对 $\omega^{-1/2}$ 作图，截距为 $1/i_k$，由此获得各过电位下的动力学电流。

\subsection{Tafel 分析}

对于单电子转移的准可逆体系，阳极 Tafel 方程为：

\begin{equation}
\eta = a + b_a \cdot \log|i_k|, \quad b_a = \frac{2.303 RT}{\alpha F} \approx \frac{0.059}{\alpha}\ \text{V/dec}
\label{eq:tafel}
\end{equation}

由 Tafel 斜率 $b_a$ 可得阳极传递系数 $\alpha$，外推 $\eta \rightarrow 0$ 可得交换电流 $i_0$。

\subsection{物理参数}

\begin{table}[h]
\centering
\caption{RDE 分析物理参数 (1M KCl, 25\textdegree C)}
\begin{tabular}{cccc}
\toprule
符号 & 参数 & 数值 & 单位 \\
\midrule
$\nu$ & 动力学粘度 & $8.9 \times 10^{-3}$ & $\text{cm}^2/\text{s}$ \\
$F$ & Faraday 常数 & $96485$ & $\text{C}/\text{mol}$ \\
$n$ & 电子转移数 & 1 & — \\
$c^*$ (Part A) & Fe(CN)$_6^{3-/4-}$ 浓度 & $1.0 \times 10^{-5}$ & $\text{mol}/\text{cm}^3$ \\
\bottomrule
\end{tabular}
\end{table}

\newpage

% ============================================================
\section{结果与讨论}

\subsection{Tafel 极化曲线特征}

0.01M K$_3$Fe(CN)$_6$/K$_4$Fe(CN)$_6$ + 1M KCl 溶液在各转速下的极化曲线呈现
典型可逆氧化还原电对特征：低电位端为阴极极限电流平台（Fe(III) $\rightarrow$ Fe(II)），
高电位端为阳极极限电流平台（Fe(II) $\rightarrow$ Fe(III)），
过渡区电流随电位快速变化。平衡电位 $E_0 \approx 0.245$ V（vs Hg/HgO），
随转速增大极限电流按 Levich 关系递增。

\begin{figure}[htbp]
\centering
\includegraphics[width=0.85\textwidth]{Fig0_Polarization_1组.jpg}
\caption{Grp1 — 0.01M K$_3$Fe(CN)$_6$/K$_4$Fe(CN)$_6$ + 1M KCl 多转速极化曲线叠加。
蓝色/红色阴影分别为阴阳极极限电流提取区间。}
\label{fig:polarization}
\end{figure}

\subsection{Levich 分析 — 扩散系数}

\begin{table}[h]
\centering
\caption{扩散系数测定结果 (4组独立测量)}
\begin{tabular}{lcc}
\toprule
操作者 & $D_{\text{O}}$ ($\times 10^{-6}$ cm$^2$/s) & $D_{\text{R}}$ ($\times 10^{-6}$ cm$^2$/s) \\
\midrule
Grp1 & 6.56 & 4.16 \\
Grp2 & 6.53 & 4.15 \\
Grp3 & 6.55 & 4.15 \\
Grp4 & 6.54 & 4.15 \\
\midrule
\textbf{平均值} & \textbf{6.55} & \textbf{4.15} \\
文献值\cite{bard2001} & 7.6 & 6.3 \\
\bottomrule
\end{tabular}
\label{tab:diffusion}
\end{table}

Levich 图（图~\ref{fig:levich}）的线性拟合 $R^2 \approx 1.0$，
表明 $|i_l| \propto \omega^{1/2}$ 关系严格成立，验证了 Levich 方程的适用性。

$D_{\text{O}}$ 比文献值低约 14\%，偏差在合理范围（温度、浓度、电极面积等因素）。
$D_{\text{R}}$ 比文献值低约 34\%，可能原因：（1）阳极极限电流平台在 0.5 V 处
未完全达到稳态；（2）Fe(CN)$_6^{4-}$ 的氧化产物可能在电极表面累积；
（3）未扣除 KCl 基底的背景电流（约 $\mu$A 级）对低转速数据影响更大。

\begin{figure}[htbp]
\centering
\includegraphics[width=0.95\textwidth]{Fig1_Levich.jpg}
\caption{Levich 图：阴极 $|i_l|$ vs $\omega^{1/2}$（左）与阳极 $|i_l|$ vs $\omega^{1/2}$（右）。
4 条线几乎重合，反映 4 组独立测量的高度一致性。}
\label{fig:levich}
\end{figure}

\subsection{浓度校准}

不同浓度（0.0002--0.01M）在 2000 rpm 下的极限电流与浓度呈良好线性关系（图~\ref{fig:conc}），
验证了 Levich 方程中 $i_l \propto c^*$ 的关系。

\begin{figure}[htbp]
\centering
\includegraphics[width=0.95\textwidth]{Fig2_Concentration.jpg}
\caption{浓度校准：$|i_l|$ vs $c^*$ @2000 rpm}
\label{fig:conc}
\end{figure}

\subsection{Koutecky-Levich 分析}

在 6 个过电位（$\eta = 80$--$200$ mV）下进行 K-L 分析（图~\ref{fig:kl}），
$1/|i|$ 对 $\omega^{-1/2}$ 的线性拟合良好（$R^2 > 0.86$），
由截距获得各过电位下的动力学电流 $i_k$。

\begin{figure}[htbp]
\centering
\includegraphics[width=0.85\textwidth]{Fig3_KL_1组.jpg}
\caption{K-L 图（Grp1）：$1/|i|$ vs $\omega^{-1/2}$，6 个过电位各一条拟合线}
\label{fig:kl}
\end{figure}

\subsection{Tafel 分析}

将 K-L 获得的 $i_k$ 代入 Tafel 分析，$\eta$ vs $\log|i_k|$ 线性拟合结果如下：

\begin{table}[h]
\centering
\caption{Tafel 分析结果}
\begin{tabular}{lcccc}
\toprule
操作者 & $b$ (mV/dec) & $\alpha$ & $R^2$ & $i_0$ ($\mu$A) \\
\midrule
Grp1 & 60.2 & 0.980 & 0.985 & 945 \\
Grp2 & 81.6 & 0.723 & 0.861 & 2311 \\
Grp3 & 81.8 & 0.721 & 0.987 & 2357 \\
Grp4 & 62.2 & 0.948 & 0.997 & 995 \\
\bottomrule
\end{tabular}
\end{table}

测得 $\alpha = 0.72$--$0.98$，偏离理想可逆值 $\alpha = 0.5$（此时 $b = 118$ mV/dec）。
原因分析：（1）Fe(CN)$_6^{3-/4-}$ 在 Pt 上为快速准可逆体系，
在所选过电位范围（$\eta = 80$--$200$ mV）内反反应（还原支）仍有不可忽略的贡献，
导致表观 Tafel 斜率偏小，$\alpha$ 偏高；
（2）$i_k$ 的提取依赖于 K-L 截距，截距对测量误差敏感。
准确的 $\alpha$ 和 $i_0$ 测定宜扩大到更高过电位区或采用 Nicholson 方法\cite{bard2001}。

\begin{figure}[htbp]
\centering
\includegraphics[width=0.95\textwidth]{Fig4_Tafel_1组.jpg}
\caption{Tafel 图（Grp1）：$\eta$ vs $\log|i_k|$。虚线为 $b=118$ mV/dec ($\alpha=0.5$) 参考线。}
\label{fig:tafel}
\end{figure}

\subsection{半波电位（无支持电解质）}

0.0008M K$_3$Fe(CN)$_6$/K$_4$Fe(CN)$_6$ 在无 KCl 条件下的极化曲线
（图~\ref{fig:hw}）显示 $E_0 \approx 0.161$ V、$E_{1/2} \approx 0.172$ V。
与有 KCl 体系（$E_0 \approx 0.245$ V）相比，
平衡电位负移约 73 mV，归因于无支持电解质时溶液欧姆降增大和液接电位变化。

\begin{figure}[htbp]
\centering
\includegraphics[width=0.85\textwidth]{Fig5_HalfWave_1组.jpg}
\caption{0.0008M 无 KCl 极化曲线及半波电位}
\label{fig:hw}
\end{figure}

\subsection{4 组独立测量的可复现性}

4 组独立操作者的 Levich 分析结果（表~\ref{tab:diffusion}）高度一致
（$D_{\text{O}}$: 6.53--6.56 $\times 10^{-6}$，$D_{\text{R}}$: 4.15--4.16 $\times 10^{-6}$），
Tafel 斜率中 2 组（Grp1/Grp4）与另 2 组（Grp2/Grp3）存在系统性差异
（$60$--$62$ vs $81$--$82$ mV/dec）。
需要指出，CHI660E 以 4 位有效数字导出数据，
限制了在相同实验条件下区分各组独立测量值的精度。
组间差异尤其在 1000 rpm 处的读数不同导致了 K-L 斜率和 Tafel 参数的分化。

\newpage

% ============================================================
\section{结论}

\begin{enumerate}[leftmargin=1.5cm]
  \item 通过 Levich 分析测得 Fe(CN)$_6^{3-}$ 和 Fe(CN)$_6^{4-}$ 的扩散系数分别为
  $D_{\text{O}} = 6.55 \times 10^{-6}$ cm$^2$/s 和
  $D_{\text{R}} = 4.15 \times 10^{-6}$ cm$^2$/s，与文献值量级一致。
  \item Levich 图和浓度校准验证了 $|i_l| \propto \omega^{1/2}$ 和 $|i_l| \propto c^*$ 关系。
  \item K-L 分析结合 Tafel 外推获得的传递系数 $\alpha \approx 0.72$--$0.98$
  偏离理想可逆值，分析归因于准可逆体系在低 $\eta$ 区的反反应贡献。
  \item 4 组独立测量在 Levich 分析中高度一致，验证了实验方案的可复现性。
\end{enumerate}

\newpage

% ============================================================
\section{局限性与改进方向}

\subsection{已实现的可复现性}

4 组独立操作者遵照相同实验方案，Levich 分析得到的扩散系数差异 $< 0.5\%$，
表明 RDE 方法在跨操作者条件下具有良好的可转移性。

\subsection{当前局限}

\begin{enumerate}[leftmargin=1.5cm]
  \item \textbf{阳极平台不完整}：$D_{\text{R}}$ 系统偏低 34\%，阳极极限电流在 0.5V 处
  可能未达到完全稳态平台
  \item \textbf{无背景扣除}：KCl 基底的背景电流（$\mu$A 级）未从数据中扣除
  \item \textbf{$\alpha$ 偏差}：Tafel 分析所选的过电位范围（80--200 mV）
  对快速准可逆体系偏低，反反应贡献导致表观 $\alpha$ 偏高
  \item \textbf{CHI660E 导出精度}：4 位有效数字限制了在多数转速下区分
  各组独立测量值的能力
\end{enumerate}

\subsection{建议改进}

\begin{enumerate}[leftmargin=1.5cm]
  \item 扩大阳极扫描范围至更高电位，确保极限电流完全达到平台
  \item 扣除 KCl 基底背景电流
  \item 将 Tafel 分析的 $\eta$ 范围扩展至 200--400 mV
  \item 采用 Nicholson 方法（CV 峰电位差法）独立验证 $\alpha$ 值
\end{enumerate}

\newpage

% ============================================================
\begin{thebibliography}{5}

\bibitem{bard2001}
A.~J.~Bard, L.~R.~Faulkner,
\textit{Electrochemical Methods: Fundamentals and Applications}, 2nd ed.,
John Wiley \& Sons, New York, 2001. (第 9 章: Rotating Disk Electrode)

\bibitem{教材}
唐安平 主编,
\textit{电化学实验},
实验 26: K$_4$Fe(CN)$_6$/K$_3$Fe(CN)$_6$体系旋转圆盘电极动力学参数的测定.

\bibitem{nicholson1965}
R.~S.~Nicholson,
``Theory and Application of Cyclic Voltammetry for Measurement of Electrode Reaction Kinetics,''
\textit{Analytical Chemistry}, 1965, \textbf{37}(11), 1351--1355.

\end{thebibliography}

\newpage

% ============================================================
\appendix
\section{附录：数据分析方法论}

\subsection{自动化数据处理流水线}

本项目数据处理由 Python 脚本自动化完成：

\begin{center}
\small
\begin{tabular}{c@{\hspace{0.5cm}}c@{\hspace{0.5cm}}c@{\hspace{0.5cm}}c}
\toprule
\textbf{步骤} & \textbf{脚本} & \textbf{输入} & \textbf{输出} \\
\midrule
1. 解析 & \texttt{parse\_chi660e.py} & Tafel Plot TXT & \texttt{all\_data.csv} \\
2. Levich & \texttt{levich\_analysis.py} & \texttt{all\_data.csv} & \texttt{limiting\_currents.csv} \\
3. K-L/Tafel & \texttt{kl\_analysis.py} & \texttt{all\_data.csv} & \texttt{kl\_analysis.csv}, \texttt{tafel\_results.csv} \\
4. 图表 & \texttt{generate\_figures.py} & 各CSV & PNG 图 \\
\bottomrule
\end{tabular}
\end{center}

\vspace{0.5cm}

\textbf{关键技术特征}：

\begin{enumerate}[leftmargin=1.5cm]
  \item \textbf{参数化配置}：实验参数存储于 \texttt{params.yaml}，脚本通过
  \texttt{yaml.safe\_load()} 读取
  \item \textbf{数据格式}：中间数据全部采用 CSV 格式（可读、可查、跨平台兼容）
  \item \textbf{类型注解}：函数接口添加 Python 类型注解
  \item \textbf{单元测试}：\texttt{tests/} 目录包含 Levich 常数计算和扩散系数
  拟合的自动化验证
  \item \textbf{兼容性}：Tafel Plot 3 列格式 + 自动跳过头部参数行
\end{enumerate}

\subsection{项目结构}

\begin{verbatim}
rde-kinetics/
|-- params.yaml
|-- run_all.sh
|-- scripts/
|   |-- parse_chi660e.py       # Tafel Plot 解析
|   |-- levich_analysis.py     # Levich 分析
|   |-- kl_analysis.py         # K-L + Tafel
|   `-- generate_figures.py    # 图表生成
|-- tests/
|   `-- test_levich.py
|-- output/
|   |-- processed/  tables/  figures/
`-- report/
    `-- RDE_分析报告.tex
\end{verbatim}

\end{document}
